\OlympFirstPage

\problem{Солянка}{10.0}

\subproblem{Фрикционная сварка}{3.5}

Рассмотрим простую модель \textit{фрикционной сварки}. Имеются два одинаковых достаточно длинных стальных цилиндра радиусом $R=\qty1{\cm}$ с совпадающими осями. Один из цилиндров раскручивают с некоторой угловой скоростью $\omega$, и его придавливают торец к торцу зафиксированного цилиндра с постоянным давлением $p=\qty{20}{\MPa}$. Определите, при каком предельном $\omega$ сталь начнёт плавиться в стационарном режиме (спустя достаточно долгое время вращения цилиндра). Считайте, что температура по радиусу цилиндра быстро выравнивается. Температура окружающего воздуха $T_0=\qty{20}{\celsius}$. Для нержавеющей стали: коэффициент трения стали о сталь $\mu=0.4$, коэффициент теплопроводности $\kappa=\qty{20}{\WmK}$, коэффициент теплоотдачи в воздух при принудительной конвекции $\alpha=\qty{200}{\W\per\metre\squared\per\kelvin}$, а точка плавления $T^*=\qty{1520}{\celsius}$. Считайте параметры стали независимыми от её температуры; теплоёмкостью цилиндров пренебречь.

\textbf{Математическая подсказка:} Решение дифференциального уравнения вида $x''(t) = \lambda^2 x(t)$ для некоторой вещественной постоянной $\lambda$ имеет вид
$$
x(t) = Ae^{\lambda t} + Be^{-\lambda t}
$$
для некоторых констант интегрирования $A$ и $B$, а $e=2.718...$ --- число Эйлера.

\subproblem{Переменный (?) ток}{3.5}

\pic7R{0pt}{0.3}{figures/1b.1}

В цепи, показанной на рисунке справа, последовательно соединённые конденсатор ёмкостью $C=\qty{4.7}{\mF}$ и катушка индуктивностью $L=\qty{5.4}{\H}$ подключены к ключу $K$, который в обоих пунктах подключается к источнику постоянного напряжения $U_0=\qty{10}{\V}$ на очень короткое время $\tau=\qty5{\ms}$, затем в течение времени $T=\qty{250}{\ms}$ (считайте, что $T\gg\tau$) переключен на короткое замыкание, т.е. крайне левое вертикальное положение, а потом снова подключается к источнику, и так далее. В момент $t=0$ тока в цепи не было в обоих пунктах.

\newcounter{taskstar}
\setcounter{taskstar}{0}
\begin{list}{\textbf{(\alph{taskstar})}\ }{
    \usecounter{taskstar}
    \setlength{\leftmargin}{18pt}%
      \setlength{\labelwidth}{0pt}%
      \setlength{\labelsep}{0pt}%
      \setlength{\itemsep}{6pt}%
      \setlength{\topsep}{6pt}}
    \item Пусть конденсатор был разряжен, а ключ находился в крайнем левом вертикальном положении. В момент $t=0$ ключ замкнулся на $U_0$ и продолжил переключаться так, как описано выше. Каким станет энергия системы в момент $t=3T+\tau$ (т.е. $t=3T$ сразу после переключения $K$ на крайнее левое положение)? Чему был максимальный заряд конденсатора $q_0$ за промежуток $0\leq t \leq 3T+\tau$?
    \item Допустим, в момент $t=0$ на конденсаторе был неизвестный заряд $q_*>0$ (полярность указана на рисунке), а ключ неподвижно находился в крайнем левом вертикальном положении. Затем, в некоторый момент $t=t_*>0$, ключ замкнули на источник $U_0$ и продолжили переключать как описано выше. Оказалось, полная энергия системы осталась неизменной при переключении ключа $K$. При каких $q_*$ и $t_*$ такое возможно?
\end{list}

\subproblem{Интерференция звука}{3.0}

Две бесконечные жесткие плоскости (стены) пересекаются под прямым углом. Из бесконечности падает плоская звуковая волна, волновой вектор которой направлен точно по биссектрисе угла между плоскостями. Амплитуда скорости молекул воздуха в падающей волне равна $v$, а длина волны составляет $\lambda$. Стены идеально отражающие (отражения происходят без потерь энергии); дифракционными эффектами на расстоянии порядка $\lambda$ пренебречь.

\setcounter{taskstar}{0}
\begin{list}{\textbf{(\alph{taskstar})}\ }{
    \usecounter{taskstar}
    \setlength{\leftmargin}{18pt}%
      \setlength{\labelwidth}{0pt}%
      \setlength{\labelsep}{0pt}%
      \setlength{\itemsep}{6pt}%
      \setlength{\topsep}{6pt}}
    \item Найдите максимально возможную скорость движения воздуха в данной системе.
    \item В каких точках пространства достигается эта максимальная скорость?
\end{list}